% =============================================================================
% SECTION 6: PRODUCTION DEPLOYMENT & MONITORING MLOPS
% =============================================================================
\section{Production Deployment \& Monitoring MLOps}

% ---------------- System Architecture ----------------
\begin{frame}{Production System Architecture}
    \begin{block}{Platform Components}
        An end-to-end MLOps ecosystem designed for dynamic feature extraction, hyperparameter shifting, and localized storage.
    \end{block}
    \begin{itemize}\scriptsize
        \item \textbf{Data Ingestion Pipeline:} Handles modular ingestion via structured CSV files securely bounded by project root configurations.
        \item \textbf{Feature Engineering Pipeline:} \texttt{DynamicFeatureEngineer} transforms raw inputs into interaction terms and indicators.
        \item \textbf{Adaptive ML Optimization:} Runs sequential multi-stage Optuna tuning loops with real-time MLflow Tracking metrics.
        \item \textbf{Model Registry:} Tracks parameters and logs serialized cloudpickle models to a dedicated SQLite storage layer.
        \item \textbf{API Inference Service:} Implements microservices using FastAPI to deliver endpoints for execution and runtime statistics.
    \end{itemize}
\end{frame}

% ---------------- Prediction Workflow ----------------
\begin{frame}{Churn Prediction Workflow}
    \begin{alertblock}{Data Lifecycle Execution}
        From raw transaction parameters to calculated strategic customer retention campaigns.
    \end{alertblock}
    \begin{enumerate}\scriptsize
        \item \textbf{Customer Input Data:} Raw structured vectors pass through validation checks upon reaching API micro-routers.
        \item \textbf{Feature Transformation:} Preprocessors construct feature names out, scaling variables using target encodings.
        \item \textbf{Pipeline Type Defense:} An implicit \texttt{XGBNumericCaster} layer forces features into \texttt{float32} shapes to shield the base boosting trees.
        \item \textbf{Model Inference:} High-velocity tree tracking models run predictive passes to assign probability weights.
        \item \textbf{X-Ray Explanations:} Triggers localized tree-path dependent SHAP routines to serialize feature importances back to MLflow.
    \end{enumerate}
\end{frame}

% ---------------- Multi-Stage Training Engine ----------------
\begin{frame}{Adaptive Multi-Stage Optimization Loop}
    \begin{itemize}
        \item Implements sequential search boundary adjustment loops over 4 key phases to scale parameter ranges dynamically:
    \end{itemize}
    \begin{table}\scriptsize
        \centering
        \begin{tabular}{lcl}
            \toprule
            \textbf{Optimization Phase} & \textbf{Trials Run} & \textbf{Search Boundary Strategy} \\
            \midrule
            Base Phase & 200 & Explores baseline hyperparam defaults \\
            Phase 2 Step & 100 & Quantile profiling and range movement \\
            Phase 3 Step & 50 & Concentrates boundaries on optimum zones \\
            Phase 4 Step & 25 & Fine-tunes localized high-perf vectors \\
            \bottomrule
        \end{tabular}
        \caption{Sequential Parameter Convergence Lifecycle}
    \end{table}
    \begin{block}{Model Promotion \& Champion Guardrails}
        The tracking server queries all parent runs inside the SQLite environment. The platform automatically promotes new models only if the calculated test PR-AUC surpasses the highest historical metric.
    \end{block}
\end{frame}

% ---------------- Core Interface Routing ----------------
\begin{frame}[fragile]{Microservice Interface Endpoints}
    The delivery layer exposes clean REST schemas managed by Uvicorn.
    \scriptsize
    \begin{lstlisting}[language=Python]
# Ingestion and Retraining 
@app.post("/train")
async def train(file: UploadFile, experiment_name: str):
    # Runs ingestion, engineering, and multi-stage loops

# Real-Time Operational Predictor
@app.post("/predict")
async def predict(file: UploadFile, run_id: str):
    # Transforms arrays and produces XGBoost probabilities

# Analytics and Audit
@app.get("/experiments/best_runs")
def get_best_runs():
    # Returns history tracking logs via 
    # BestRunResponse Pydantic schema
    \end{lstlisting}
\end{frame}

% ---------------- Containerized Deployment ----------------
\begin{frame}{Containerized Deployment Architecture}
    \begin{block}{Reproducible Ecosystem Environment}
        The application is packaged inside a lightweight \texttt{python:3.12-slim} container, ensuring quick start-up times and isolation.
    \end{block}
    \begin{itemize}\scriptsize
        \item \textbf{FastAPI Application Container (\texttt{hop\_api}):} Exposes internal operational processing functions via port \texttt{8000}.
        \item \textbf{MLflow Dashboard Service (\texttt{hop\_mlflow}):} Mounts unified workspace directory binds to expose the visualization engine UI on port \texttt{5000}.
        \item \textbf{Persistent Storage Layer:} Mounts localized shared system volumes to maintain tracking information inside \texttt{mlflow.db} and model plots under \texttt{mlartifacts/}.
    \end{itemize}
\end{frame}

% ---------------- Production Monitoring ----------------
\begin{frame}{Production Monitoring}
    \begin{block}{Automated Drift Detection Framework}
        Production monitoring services detect data and target distribution shifts before model performance degradation becomes significant.
    \end{block}
    \begin{itemize}\scriptsize
        \item \textbf{Feature Drift Detection:} KS tests for numerical variables and Chi-Square tests for categorical variables.
        \item \textbf{Pipeline-Aware Monitoring:} Drift evaluated after feature engineering and preprocessing transformations.
        \item \textbf{Target Drift Analysis:} Monitoring of churn-rate shifts and class distribution changes.
        \item \textbf{Monitoring APIs:} Real-time drift assessment endpoints returning statistical reports and drift alerts.
    \end{itemize}
\end{frame}

% ---------------- Monitoring Interface Routing ----------------
\begin{frame}[fragile]{Monitoring Service Endpoints}
    The observability layer exposes statistical drift monitoring services for production model governance.
    \scriptsize
    \begin{lstlisting}[language=Python]
# Feature Distribution Monitoring
@app.post("/monitor/feature-drift")
async def monitor_drift(reference_file: UploadFile, 
                        current_file: UploadFile, run_id: str):
    # Executes KS and Chi-Square drift tests
    # across transformed production features

# Target Distribution Monitoring
@app.post("/monitor/target-drift")
async def monitor_target_drift(reference_file: UploadFile, 
                                current_file: UploadFile, 
                                target_column: str = "Exited"):
    # Evaluates churn-rate shifts and
    # class distribution drift statistics
    \end{lstlisting}
\end{frame}

% ---------------- End-to-End MLOps Pipeline ----------------
\begin{frame}{End-to-End MLOps Pipeline}
    \begin{exampleblock}{Continuous Integration \& Training Loop}
        Unified architectural lifecycle mapping feedback dependencies from collection to iteration.
    \end{exampleblock}
    \vspace{0.3cm}
    \begin{center}
        \small
        \colorbox{gray!15}{\textbf{\strut Data}} $\rightarrow$ 
        \colorbox{gray!15}{\textbf{\strut Features}} $\rightarrow$ 
        \colorbox{gray!15}{\textbf{\strut Training}} $\rightarrow$ 
        \colorbox{gray!15}{\textbf{\strut MLflow}} $\rightarrow$ 
        \colorbox{gray!15}{\textbf{\strut Deployment}} $\rightarrow$ 
        \colorbox{gray!15}{\textbf{\strut Monitoring}} $\rightarrow$ 
        \colorbox{gray!15}{\textbf{\strut Retraining}}
    \end{center}
\end{frame}